%/*% pdflatex -file-line-error -interaction nonstopmode %

\documentclass{article}
\usepackage{graphicx}
\graphicspath{ {.} }
\usepackage{amsmath}
\usepackage{amssymb}
\usepackage{amsfonts}
\usepackage{hyperref}
\usepackage{exercises}
\usepackage{amsthm}
\usepackage{yhmath}
\usepackage{placeins}
\newcommand*{\oldneg}{\mathord{\sim}}
\let\oldemptyset\emptyset

\theoremstyle{definition}
\newtheorem{definition}{Definition}

\setlength{\parindent}{.5in}
\setlength{\parskip}{12pt}
\title{Chapter 1 Answers}
\begin{document}

\section*{1}
(a): First, we show that if $S$ is a boolean truth set, then (1) and (2) follow.\newline 
(1): We are to show that $\alpha\in S$ iff $\alpha_1\in S$ and $\alpha_2\in S$. The case where $\alpha$ is $X\wedge Y$ is just condition (1). By (4), 
$\oldneg (X\supset Y)$ iff $X\supset Y\notin S$. So, $\oldneg (X\supset Y)\in X$ iff $X\supset Y\notin S$ iff $X\in S$ and $Y\not in S$ (by (3)) iff 
$\oldneg Y\in S$, which proves the case of $X\supset Y$. In the third case, the formula is of the form $\oldneg (X\vee Y)$. $\oldneg (X\vee Y)\in S$ iff 
$X\vee Y\notin S$ (by (4)) iff $X\notin S$ and $Y\notin S$ (by (2)) iff $\oldneg X\in S$ and $\oldneg Y\in S$ (by (4)), which completes this case. In the 
last case, the formula is $\oldneg\oldneg X$. $\oldneg\oldneg X\in S$ iff $\oldneg X\notin S$ (by (4)) iff $X\in S$ (by (4)), which completes the proof in 
this case.\newline
(2): We are to show that either $X\in S$ or $(\oldneg X)\in S$, but not both. By (4), $(\oldneg X)\in S$ iff $X\notin S$. 
In any case, this shows that either $X\in S$ and $(\oldneg X)\notin S$, or that $(\oldneg X)\in S$ and $X\notin S$ (by definition of iff), which is 
what we were trying to prove.
(b): Next, suppose that the two conditions of Problem 1 hold. We are to show that $S$ is a Boolean truth set. 
(1): This is immediate from condition (1) of Problem 1.\newline
(2): By (2) from Problem 1 $X\vee Y\in S$ iff $\oldneg(X\vee Y)\notin S$ iff $\oldneg X\notin S$ or $\oldneg Y\notin S$ (by (1) of Problem 1) iff 
$X\in S$ or $Y\in S$ (by (2) of Problem 1)\newline
(3): By (2) from Problem 1 $X\supset Y\in S$ iff $\oldneg(X\supset Y)\notin S$ iff $X\notin S$ or $\oldneg Y\notin S$ (by (1) of Problem 1) iff 
$\oldneg X\in S$ of $Y\in S$ (by (2) of Problem 1)\newline
(4): This is obviously equivalent to (4).

\section*{Exercise 1}
(a): First, we show that if $S$ is a boolean truth set, then (1) and (2) from this exercise follow.\newline
(1): The cases where $\beta$ is of type $X\vee Y$ or $X\supset Y$ are immediate (by (2) and (3) from the Boolean consequence relation definition). 
When $\beta$ is of the form $\oldneg(X\wedge Y)$, $\oldneg(X\wedge Y)\in S$ iff $X\wedge Y\notin S$ (by (4)) iff $X\notin S$ or $Y\notin S$ (by (1)) 
iff $\oldneg X\in S$ or $\oldneg Y\in S$ (by (4)).\newline
(2): This is equivalent to (4) in the definition of a Boolean consequence relation.\newline
(b): Next, we show that if $S$ satisfies (1) and (2) from this exercise, then $S$ is a Boolean truth set.
(1): $X\wedge Y\in S$ iff $\oldneg(X\wedge Y)\notin S$ (by (2)) iff $\oldneg X\notin S$ and $\oldneg Y\notin S$ (by (1)) iff $X\in S$ and $Y\in S$ (by 
(2)).\newline
(2-3): Conditions (2) and (3) are immediate from condition (1) of this exercise. 
(4): This is equivalent to condition (2) of this exercise.

\section*{2}
\[
\begin{array}{rl}
	& 1:\ S\vdash X\ (Hypothesis) \\
	& 2:\ S\vdash Y\ (Hypothesis) \\
	& 3:\ X,Y\vdash Z\ (Hypothesis) \\
	& 4:\ S,X,Y\vdash X,Y\ (C_1) \\
	& 5:\ S,X\vdash Y\ (2,C_2) \\
	& 6:\ S,X\vdash X,Y\ (5,4,C_3) \\
	& 7:\ S\vdash X,Y\ (1,6,C_3) \\
	& 8:\ S,X,Y\vdash Z\ (3,C_2) \\
	& 9:\ S\vdash Z\ (7,8,C_3)
\end{array}
\]

\section*{3}
We assume that $S$ is consistent, and $S\vdash X$. We must show that $S,X$ is consistent. Because $S$ is consistent, this means that $S\vdash Z$ is not the 
for at least one formula $Z$. Now suppose that $S,X\vdash Z$ is the case. By $C_3$, this would mean that $S\vdash Z$, which is a contradiction. This means 
that $S,X\vdash Z$ is not the case, meaning $S,X$ does not yield at least one formula $Z$, thus proving that it is consistent.

\section*{4}
(a): We assume that S is consistent,
$S\vdash X$ and $X\vdash Y$. By $C_2$, $S,X\vdash Y$ follows. By $C_3$, $S\vdash Y$ also follows, meaning $S,Y$ is consistent by 
Problem 3.\newline
(b): We assume that S is consistent,  
$S\vdash X$ and $S\vdash Y$ and $X,Y\vdash Z$. By Problem 2, $S\vdash Z$ follows, and thus, by Problem 3, $S,Z$ is consistent.

\section*{5}
(a): We assume that M is maximally consistent, 
$X\in M$ and $X\vdash Y$ and that $M$ is maximally consistent. Suppose that $Y\notin M$. By $C_1$, $M\vdash X$, meaning, by Problem 4(a),
that $M,Y$ is consistent. If we assume that $Y\notin M$, this is in direct contradiction with the fact that $M$ is maximally consistent, thus meaning 
$Y\in M$.\newline
(b): We assume that M is maximally consistent, $X\in M$ and $Y\in M$ and $X,Y\vdash Z$. This means that $M\vdash X$ and $M\vdash Y$ by $C_1$. By 
Problem 4(b), it follows that $M,Z$ is consistent. So, if $Z\notin M$, this is in direct contradiction with the fact that $M$ is maximally consistent, 
meaning $Z\in M$. 

\section*{Exericse 2}
We need to show that $S\vDash X$ is a Boolean consequence relation, that is, that it satisfies $C_1-C_8$.\newline
$C_1$: Suppose $X\in S$. Because $X\in S$, $X$ is obviously true under all interpretations that satisfy $S$, meaning $S\vDash X$.\newline
$C_2$: Suppose $S_1\vDash X$ and $S_1\subseteq S_2$. This means that $X$ is true under all interpretations that satisfy $S_1$. Well, because 
$S_1\subseteq S_2$, all interpretations that satisfy $S_2$ also satisfy $S_1$, meaning $X$ is also true under all interpreations that satisfy 
$S_2$, thus proving $S_2\vDash X$.\newline
$C_3$: Suppose $S\vDash X$ and $S,X\vDash Y$. This means that $Y$ is true under all interpretations that satisfy $S$ and satisfy $X$. But 
$X$ is true under all interpreations that satisfy $S$, meaning, under all interpreations that satisfy $S$, $S,X$ is also satisfied. Thus, 
$S\vDash Y$.\newline
$C_4$: Suppose $S\vDash X$. Recall the Corollary to the Compactness Theorem for Propositional Logic that states that 
if $X$ is a logical consequence of $S$ then $X$ is a logical consequence of some finite subset of $S$ (logical consequence there means the same 
thing as tautological consequence here), which single-handedly proves $C_4$.\newline
$C_5$: (a-b): By definition of an $\alpha$ formula, it is true iff $\alpha_1,\alpha_2$ are true, meaning that $\alpha\vDash\alpha_1,\alpha\vDash\alpha_2$.
\newline
$C_6$: Again, by definition of an $\alpha$ formula, it is true iff $\alpha_1,\alpha_2$ are true, meaning $\alpha_1,\alpha_2\vDash\alpha$.\newline
$C_7$: By definition, $X,\oldneg X$ is unsatisfiable. Hence, every formula $Y$ is vacuously true under all interpreations that satisfy $X,\oldneg X$ 
(because there aren't any). \newline
$C_8$: Suppose that $S,\oldneg X\vDash X$. This means that, if $\oldneg X$ was true under any interpretations that satisfied $S$, $X$ would also be 
true, yielding a contradiction. Hence, $S\bigcup\oldneg X$ is unsatisfiable, meaning $X$ is true under all interpretations that satisfy $S$.

\section*{6}
We suppose that $\vdash$ is tautologically complete. This means that, whenver $X$ is a tautological consequence of $\oldemptyset$. Because we are considering 
the empty set, this means that all elements of the empty set of vacuously satisfied by every interpretation. Because $X$ is assumed to be a tautology, 
it to is satisfied by every interpretation, meaning $X$ is a tautological consequence of $\empty$, and thus, because $\vdash$ is tautologically complete, 
$\oldemptyset\vdash X$.

\section*{7}
We must prove that if $\vdash$ is a Boolean consequence relation and $M$ is maximally consistent with respect to $\vdash$, then $M$ is a Boolean truth 
set. Recall that because $\vdash$ is a Boolean consequence realtion, it satisfies $C_1-C_8$. Also, because $M$ is maximally consistent with respect to 
$\vdash$, it is not a proper subset of any consistent set with respect to $\vdash$ such that this set has an element not in $M$. We need to show that 
$M$ is a Boolean truth set, which, by Problem 1, is the same as showing that it satisfies conditions (1) and (2) of that problem. As for condition (1), 
we need to show that $\alpha\in S$ iff $a_1\in S$ and $a_2\in S$. On the one hand, assume that $\alpha\in M$. By $C_1$, $M\vdash\alpha$. By $C_5$, 
$\alpha\vdash\alpha_1$ and $\alpha\vdash\alpha_2$. Hence, by two applications of Problem 5(a), we can conclude that $\alpha_1\in M$ and $\alpha_2\in M$. 
On the other hand, assume that $\alpha_1\in M$ and $\alpha_2\in M$. Again, by $C_1$, $M\vdash\alpha_1,M\vdash\alpha_2$. 
By $C_6$, $\alpha_1,\alpha_2\vdash\alpha$. 
Then, by Problem 5(b), we can conclude that $\alpha\in M$.\newline
Now, we need to show that $M$ satisfies condition (2) of Problem 1. This condition states that either $X\in S$ or $(\oldneg X)\in S$, but not both. 
Suppose both $X$ and $\oldneg X$ were in $M$. By $C_7$, $X,\oldneg X\vdash Y$. Because $M$ is maximally consistent, it is consistent (with repsect to 
$\vdash$), meaning $M\nvdash Z$ for some formula $Z$. However, by $C_7$, $X,\oldneg X\vdash Z$, meaning, if $X,\oldneg X\in M$, then 
$M\vdash Z$ by $C_2$, which is a contradiction.\newline
Next, we need to show that one and only one of $X,\oldneg X$ are in $M$. First, assume that $\oldneg X\notin M$. We need to show that $X\in M$. 
Because $M$ is maximally consistent, $M,\oldneg X$ must be inconsistent (because $\oldneg X\notin M$). Hence, $M,\oldneg X\vdash Y$ holds for 
every $Y$, and in particular, $M,\oldneg X\vdash X$ holds, meaning, by $C_8$, $M\vdash X$ holds. By Problem 5(a), substituting $M$ for $X$, 
it follows that $X\in M$. The proof for the case when $X\notin M$ is similar, replacing $X$ by $\oldneg X$, and $\oldneg X$ by $X$.\newline
$Alternatively$: To show that $M\vdash X$ means $X\in M$, one can also use $C_3$. In particular, recall that $M\nvdash Z$ for the $Z$ defined above. 
Then, if $X\notin M$, this means $M,X$ must be inconsistent (because $M$ is maximally consistent), and thus $M,X\vdash Z$. By $C_3$, it would follow 
that $M\vdash Z$, which is a contradiction, meaning $X\in M$. 

\section*{8}
We are to show that for a Boolean consequence relation $\vdash$, consistency is compact, that is, a set $S$ is consistent (with respect to $\vdash$) iff 
all finite subsets of $S$ are consistent. On the one hand, suppose $S$ is consistent. Thus, $S\nvdash Z$ for some $Z$. Suppose at least one finite 
subset $F$ of $S$ is inconsistent, that is, is such that $F\vdash Y$ is provable for every $Y$. In particular, $F\vdash Z$ holds, meaning, by 
$C_2$, $S\vdash Z$ holds, which is a contradiction. Thus, all finite subsets of $S$ are consistent. On the other hand, suppose all finite subsets of 
$S$ are consistent. Suppose that $S$ is inconsistent, that is, $S\vdash Y$ holds for every formula $Y$. Because all finite subsets are consistent, 
$F\nvdash Z$ for every finite subset $Z$ for some $Z$ (possibly a different $Z$) for every finite subset.  

\end{document}




















